\scanpage{36}
\mspage{17}
peut être préférable -- l'an dernier j'ai choisi délibérément l'image de $\deleted{\unsure{\mathrm R}}\rho_\infty$ comme \inserted{$\lambda_\infty=$} $\rho^{-1}$ et non $-\rho^{-1}$, pour avoir un élément d'ordre 6, ``moins spécial'' donc que $-\lambda_\infty$ qui est d'ordre 3.

Relativement au choix (87)\inserted{$=(69)$} pour les $\rho_i$, on a le choix correspondant \inserted{$(87\ \mathrm{bis})=$}(72) pour les $\tau_i$, qui semble plus satisfaisant que l'autre choix possible
\[
\tau_i\mapsto-r_i,
\]
qui introduit plus de signes dans les expressions matricielles. Relativement au choix $(88)=(83)$, \unsure{le} choix \deleted{\unsure{cor}} compatible pour les $\tau_i$ est donné par (82), \deleted{\unsure{qu'on p|qui}}
\[
\tau_0\mapsto r_0,\quad\tau_1\longrightarrow-\unsure{\tau_1|r_1},\quad\tau_\infty\mapsto-r_\infty
\]
qu'on pourrait aussi remplacer par
\[
\tau_0\mapsto\unsure{-|\deleted{-}}r_0,\quad\tau_1\longrightarrow r_1,\quad\tau_\infty\mapsto r_\infty,
\]
mais qui n'a pas une signification géométrique aussi simple, en termes d'une action sur un $(x,y,z)\in B\deleted{\unsure{_{\varphi}}}(M)$\,; comme étant resp. la transposition de $x$ et $y$, de $y$ et $z$, ou la symétrie p.r. à l'axe des $x$

\deleted{(symétries p.r. à la diagonale principale, la diagonale non principale, et l'axe des $x$ dans le syst. de coordonnées défini par la base $(x,y)$ de $M$ \ldots)}
% archive sheet number: 35
